\chapter{Conclusiones y Recomendaciones}
\label{c5} % la etiqueta para referencias


La presente investigación logró cumplir con los objetivos propuestos, al desarrollar una formulación que permite representar de manera más realista el comportamiento de un planificador social dentro de un sistema de permisos de emisión transables. En particular, el trabajo incorpora la teoría de inatención conductual al proceso de fijación del límite de emisiones, reemplazando el supuesto de información perfecta presente en formulaciones previas. Bajo este enfoque de \textit{anclaje y ajuste}, el regulador parte desde una estimación inicial del presupuesto de carbono natural $CAP$ y ajusta su decisión en función del nivel de atención que dedica al procesamiento de la información disponible. De esta manera, el modelo muestra que la precisión de la meta ambiental depende directamente del grado de análisis que el regulador realiza sobre las señales informativas del sistema. Los resultados indican que mayores niveles de atención permiten reducir la incertidumbre en la estimación del presupuesto de emisiones, lo que se traduce en una menor participación de tecnologías contaminantes y una mayor presencia de tecnologías limpias en la matriz energética. \medskip

Al calibrar el modelo con datos de la matriz productiva de Chile al año 2024, los resultados sugieren que actualmente existirían niveles bajos de atención hacia las señales asociadas al presupuesto real de emisiones. En particular, el modelo estima un presupuesto de 311,01 MtCO$_2$e para el período analizado, valor que supera el límite de 264,2 MtCO$_2$e establecido por el \citeA{ministerio_de_energia_de_chile_plan_2024} para el sector eléctrico. Esta brecha sugiere que la dinámica actual del sistema energético se mantiene anclada en trayectorias históricas de emisiones, sin converger plenamente hacia los objetivos ambientales definidos por la política climática del país. En este contexto, los resultados también sugieren que instrumentos basados únicamente en impuestos a las emisiones pueden resultar insuficientes para inducir cambios significativos en la matriz energética, particularmente cuando dichos costos pueden ser parcialmente trasladados a los consumidores finales. Frente a ello, mecanismos de regulación basados en límites cuantitativos, como los sistemas de \textit{Cap-and-Trade}, presentan la ventaja de fijar un techo de emisiones no flexible que permite asegurar el cumplimiento de las metas ambientales establecidas. \medskip

Finalmente, como línea futura de investigación, sería relevante extender el horizonte temporal del modelo hacia el período 2025--2050 con el objetivo de analizar qué niveles de atención permitirían alcanzar la carbono-neutralidad. Este análisis permitiría evaluar si existe un umbral específico de atención necesario para cumplir dicha meta o si el objetivo puede alcanzarse dentro de un rango determinado de niveles de atención. Asimismo, otra extensión de este trabajo podría ser la de incorporar explícitamente el costo de aumentar la capacidad de atención del planificador, lo que permitiría cuantificar los recursos económicos necesarios para mejorar la precisión de las decisiones regulatorias y fortalecer la efectividad de la política climática.









